% ============================================================
% BAB KESIMPULAN DAN SARAN
% ============================================================
% NOTE (2026-03-04): Kesimpulan ini ditulis berdasarkan hasil
% eksperimen skala kecil (3 tanggal Januari 2022, 15 epoch).
% Bagian yang ditandai [*] perlu diperbarui setelah eksperimen
% skala penuh dengan 72 tanggal selesai.
% ============================================================

\section{Kesimpulan}

Berdasarkan hasil eksperimen dan analisis yang telah dilakukan, penelitian ini menghasilkan
kesimpulan sebagai berikut yang menjawab tiga rumusan masalah penelitian:

\begin{enumerate}

    \item \textbf{Perbandingan akurasi \textit{single-date} versus \textit{multi-temporal}.}

    Pada eksperimen skala kecil dengan tiga tanggal akuisisi Januari 2022 dan 10 kelas
    tanaman aktif (tanpa \textit{Fallow/Idle}), pendekatan \textit{multi-temporal} naif
    (Exp B, 27 channel, mIoU SegFormer = 0.2519) tidak menunjukkan keunggulan yang
    signifikan dibandingkan pendekatan \textit{single-date}
    (Exp A, 9 channel, mIoU SegFormer = 0.2556).
    Hasil ini disebabkan oleh keterbatasan variasi temporal --- ketiga tanggal yang tersedia
    semuanya berasal dari bulan Januari sehingga tidak merepresentasikan siklus fenologi
    lengkap.
    Bukti manfaat multi-temporal paling terlihat pada kelas \textit{Winter Wheat}, yang
    mengalami peningkatan IoU terbesar dari Exp A ke Exp B (27.13\% $\rightarrow$ 35.52\%
    pada SegFormer), mengonfirmasi bahwa informasi temporal yang relevan membantu
    membedakan tanaman dengan karakteristik fenologi musim yang berbeda.
    Eksperimen skala penuh dengan 72 tanggal yang mencakup seluruh siklus tanam
    diharapkan menunjukkan keunggulan multi-temporal yang lebih nyata.~[*]

    \item \textbf{Efektivitas \textit{pipeline} seleksi band tiga tahap.}

    \textit{Pipeline} seleksi band yang diusulkan --- terdiri dari \textit{Global Separation
    Index} (GSI), \textit{Random Forest feature importance}, dan \textit{CNN forward
    selection} berbasis mIoU --- berhasil mengidentifikasi subset band optimal sebanyak
    $K^* = 4$ band dari total fitur yang tersedia.
    Band yang terpilih adalah B11 (1 Jan), B12 (1 Jan), B12 (31 Jan), dan B8A (1 Jan),
    yang merepresentasikan region SWIR dan \textit{Red Edge} --- dua kawasan spektral
    paling diskriminatif untuk vegetasi pertanian.
    Dengan hanya 4 channel (efisiensi 85.2\% dari 27 channel Exp B), Exp C mempertahankan
    akurasi yang sebanding: penurunan mIoU hanya 1.98 poin pada SegFormer
    (0.2519 $\rightarrow$ 0.2321).
    Hasil ini memvalidasi hipotesis kedua bahwa subset band hasil seleksi dapat
    mempertahankan akurasi segmentasi yang kompetitif dengan beban komputasi yang
    jauh lebih rendah.

    \item \textbf{Perbandingan kinerja DeepLabV3+CBAM dan SegFormer.}

    SegFormer (encoder MiT-B2) secara konsisten mengungguli DeepLabV3+CBAM (encoder
    ResNet-50 dengan CBAM) pada seluruh konfigurasi input.
    Selisih mIoU berkisar antara 1.49 hingga 2.83 poin: Exp A (+2.83 poin),
    Exp B (+1.49 poin), dan Exp C (+1.66 poin).
    Keunggulan ini konsisten dengan hipotesis ketiga bahwa mekanisme
    \textit{Efficient Self-Attention} pada SegFormer lebih efektif menangkap konteks
    spasial global yang relevan untuk membedakan petak-petak lahan pertanian berskala
    besar dalam citra Sentinel-2.
    Sebaliknya, keunggulan konvolusi lokal DeepLabV3+CBAM relatif lebih terbatas pada
    data resolusi 10 meter yang memiliki batas antar kelas yang gradual.

\end{enumerate}

Secara keseluruhan, penelitian ini membuktikan bahwa \textit{pipeline} tiga tahap yang
mengintegrasikan seleksi fitur spektral-temporal dengan model \textit{semantic segmentation}
adalah pendekatan yang layak untuk pemetaan tanaman berbasis citra Sentinel-2.
SegFormer dengan input 4 channel hasil seleksi terbukti mampu menghasilkan akurasi
yang mendekati konfigurasi penuh dengan efisiensi komputasi yang lebih tinggi.

\section{Saran}

Berdasarkan temuan dan keterbatasan penelitian ini, berikut adalah saran untuk
pengembangan penelitian selanjutnya:

\begin{enumerate}

    \item \textbf{Eksperimen skala penuh dengan data 72 tanggal.}
    Eksperimen skala kecil ini perlu diulang menggunakan seluruh dataset multi-temporal
    (72 tanggal, 2022--2024) agar pengaruh informasi fenologi sepanjang siklus musim
    tanam dapat terukur secara definitif.
    Hasil eksperimen skala penuh akan menjadi dasar kesimpulan akhir yang lebih valid
    mengenai keunggulan pendekatan \textit{multi-temporal}.

    \item \textbf{Penanganan ketidakseimbangan kelas.}
    Kelas minoritas seperti \textit{Other Hay} (IoU 3--4\%) dan \textit{Alfalfa}
    (IoU 8--18\%) secara konsisten menunjukkan performa rendah akibat distribusi kelas
    yang tidak seimbang.
    Penelitian selanjutnya disarankan mengeksplorasi strategi seperti \textit{class-weighted
    loss}, \textit{focal loss}, atau oversampling patch berbasis kelas minoritas untuk
    meningkatkan akurasi pada kelas yang langka.

    \item \textbf{Pengujian generalisasi antar area studi.}
    Penelitian ini dibatasi pada Sacramento Valley, California.
    Untuk mengevaluasi kemampuan generalisasi model dan \textit{pipeline} seleksi band,
    perlu dilakukan pengujian pada area studi lain --- misalnya San Joaquin Valley
    (California) atau kawasan pertanian di luar Amerika Serikat --- dengan pola tanam
    dan jenis tanaman yang berbeda.

    \item \textbf{Pengembangan \textit{pipeline} seleksi band.}
    \textit{Pipeline} seleksi band yang diusulkan menggunakan U-Net ResNet-18 sebagai
    \textit{oracle} pada Stage 2 untuk efisiensi komputasi.
    Penelitian selanjutnya dapat membandingkan oracle ini dengan model alternatif
    (misalnya SegFormer-B0) atau mengganti \textit{sequential forward selection} dengan
    metode yang lebih efisien secara komputasi seperti \textit{mutual information}-based
    selection atau \textit{LASSO regularization} agar seleksi dapat dilakukan pada ruang
    fitur yang lebih besar.

    \item \textbf{Integrasi \textit{encoder} temporal eksplisit.}
    Pendekatan \textit{multi-temporal} pada penelitian ini menggunakan penggabungan
    band secara naif (\textit{channel stacking}).
    Penggunaan modul temporal eksplisit --- seperti \textit{Long Short-Term Memory} (LSTM),
    \textit{Temporal Self-Attention}, atau \textit{TimeSformer} --- yang secara arsitektural
    memisahkan dimensi temporal dan spasial diharapkan dapat mengekstraksi pola fenologi
    lebih efektif dibandingkan \textit{channel stacking}.

\end{enumerate}
